% ID: SEL-D0014-BC-E06
% ESTADO: validada. S015; tres pasadas del mismo autor, sin segundo revisor.
\textbf{Recuentos por sabor.}\par\nopagebreak[4]
Hay el mismo número de latas de cada sabor: $30000/3=10000$. Sea $E$ caducar en el año indicado en el ejercicio. Los recuentos son
\[
\begin{array}{c|rr|r}
&E&E^c&\text{Total}\\\hline
\text{Cola}&1800&8200&10000\\
\text{Naranja}&2400&7600&10000\\
\text{Limón}&3000&7000&10000\\\hline
\text{Total}&7200&22800&30000
\end{array}
\]
\textbf{a) Caducar en 2021.}\par\nopagebreak[4]
\[
\boxed{P(E)=\frac{1800+2400+3000}{30000}=\frac{7200}{30000}=0.24.}
\]
Se comprueba también por probabilidad total con pesos $1/3$ y tasas $0.18,0.24,0.30$.

\textbf{b) Cola, condicionado a no caducar en 2021.}\par\nopagebreak[4]
El nuevo espacio de referencia son las $22800$ latas que no cumplen $E$:
\[
\boxed{P(\text{cola}\mid E^c)=\frac{8200}{22800}=\frac{41}{114}\approx0.359649.}
\]
La fecha forma parte del relato original y no cambia el curso del examen ni se interpreta como dato actual de existencias.
